Considérons le polynôme $P(n) = n^2 - n + 5$.
\begin{itemize}
  \item Vérifiez que pour les premières valeurs de $n$, $P(n)$ est premier.
  \item \textbf{Discutez :} un \emph{grand} nombre d'exemples favorables
  suffirait-il à prouver que $P(n)$ est toujours premier ?
  \item Cherchez un entier $n$ rendant $P(n)$ composé.
\end{itemize}
Reprendre ces questions avec $P(n) = n^2 - n + 11$, $P(n) = 2^{2^n} + 1$.

Est-il vrai que si $p$ est un nombre premier, alors $2^p - 1$ est premier.
